\section{A compact-source obstruction even in the unrestricted strong completion}
\label{sec:r11-strong-completion-source-obstruction}

This section specializes only to the explicit counterexample background
(DO6), with $n\geq2$, $Q=(-1,1)^n$, and remaining degree $q=1$.
The field, norm, analytic space, and original test category are unchanged.
Write $U=t^{-1}(Q)$ and let $\sigma:Q\to U$ be the joint original
Gauss section.  The new potential space is precisely
$\widehat{\mathcal E}^{1}_{\rm str}(U)$ of
Definition~\ref{def:r11-strong-first-energy-completion}.
Its elements need not have an actual common lower barrier or an
ambient classical representative.  Their currents and multipliers
are those already constructed in (IC10)--(IC13).

The frozen inputs are exactly the R6 joint-Gauss convexity and
original calibration, and the R10 compact-source convex rigidity
listed in the power-class obstruction section.  The already proved
actual and fixed-barrier obstructions (DO7) and (BV13) are not
reproved here.  The new issue is whether the energy bound of an
arbitrary strongly Cauchy row still forces a finite convex trace
in its current limit.  The following estimate answers that issue.

\begin{lemma}[Local trace compactness from a first-energy bound]
\label{lem:r11-energy-bounded-convex-trace-compactness}
Let $(u_j)\subset\mathcal E^1_{S,\partial}(U)$ be actual classical
potentials with $\sup_j e_1(u_j)\leq M<\infty$.
Their finite convex Gauss traces $g_j=u_j\circ\sigma\leq0$ have
uniformly bounded $L^1$ norms on every compact real subbox of $Q$.
They are locally uniformly bounded and locally equicontinuous.
Consequently a subsequence converges locally uniformly on $Q$
to a finite convex function $g\leq0$.
No common pointwise lower barrier is assumed.
\end{lemma}

\noindent\textbf{Status: \texttt{NEW\_PROVED}} in this precise
original moment example.  The compactness assertion concerns real
convex traces, not an asserted compactness of ambient classical psh
potentials in their original category.

\begin{proof}
The finite-convex-trace argument preceding (DO8) applies separately
to every actual $u_j$: its compact negative sublevels supply one
finite Gauss value, and extended convexity bounded above by zero
then gives finiteness everywhere.  Fix a compact real subbox
$B\Subset Q$ and put $K=\sigma(B)$, a compact ordinary analytic
set.  The normalized smooth original tester
\[
                v_0=\frac1n\sum_{i=1}^n t_i^2
           \quad\hbox{satisfies}\quad
                D(v_0S)=2\sigma_*(dy).
 \tag{SO1}
\]
Let $\rho$ be the strict-box barrier from (TC2).
Choose $A>0$ so that $A\rho<-2$ near $K$ and set
\[
                     \psi=\max(v_0-2,A\rho).
 \tag{SO2}
\]
The exact core construction (TC7)--(TC8), with its harmless
constant rescaling, shows that $\psi$ is an actual bounded
finite-mass member of every positive power class.  In particular
$e_1(\psi)<\infty$, and it agrees with $v_0-2$ near $K$.
Original locality therefore gives
$\mathbf1_K\mu_\psi=2\sigma_*(\mathbf1_Bdy)$.
The global constant-one mixed energy inequality yields
\[
 \int_B(-g_j)\,dy
   =\frac12\int_K(-u_j)\,d\mu_\psi
   \leq\frac12\int_U(-u_j)\,d\mu_\psi
   \leq\frac12M^{1/2}e_1(\psi)^{1/2}.
 \tag{SO3}
\]
Every multiplier in this equation is the previously justified
actual multiplier.  This proves the local $L^1$ bound, including
the case $M=0$.

For completeness, the convex compactness step is elementary.
Fix a compact real set $C\Subset\operatorname{Int}(B)$ and
choose $\delta>0$ such that each Euclidean ball
$B(x,\delta)$ with $x\in C$ lies in $B$.
The functions $h_j=-g_j$ are nonnegative and concave.
For $z\in B(x,\delta/2)$, write $z=(x+y)/2$ with
$y\in B(x,\delta)$.  Concavity gives
$h_j(z)\geq h_j(x)/2$.  Integration implies
\[
 0\leq -g_j(x)
       \leq\frac{2}{\operatorname{vol}(B(0,\delta/2))}
                       \int_B(-g_j)\,dy.
 \tag{SO4}
\]
Thus the traces have uniform two-sided bounds on every compact
interior set.  Apply this bound on a slightly larger compact
set.  On any line through a point of the smaller set, the
monotone secant slopes of a convex function are then bounded
by its oscillation on the larger set divided by the fixed
distance to its boundary.  This gives a uniform Lipschitz
bound on every smaller compact subbox.  A finite-net
subsequence argument and a countable exhaustion yield local
uniform convergence of a subsequence.  Convexity and the
upper bound zero pass to its finite limit.
\end{proof}

\begin{theorem}[Strong completion does not create a nonzero compact source]
\label{thm:r11-strong-completed-compact-source-obstruction}
For the precise background just specified, let
$u\in\widehat{\mathcal E}^{1}_{\rm str}(U)$.
If the positive original measure $D(uS)=\widehat N_1(u)$
has tropical pushforward carried by a compact subset of $Q$,
then
\[
                              D(uS)=0.
 \tag{SO5}
\]
In particular, the nonzero positive compact source
$\mu=\sigma_*(\chi\,dy)$ of (DO9) is not $D(uS)$ for
any member of this unrestricted strong completion.
Thus passing from actual power-class potentials to the complete
relative strong first-energy space does not make finite-dual or
capacity-controlled data universally solvable.
\end{theorem}

\noindent\textbf{Status: \texttt{NEW\_PROVED}} as an extension
of the frozen source obstruction to the new strong completion
without an actual lower barrier.  The universal source-solvability
inference being rejected is \texttt{FALSE}.

\begin{proof}
Choose an actual strongly Cauchy defining row $(u_j)$ for $u$.
Its first energies are bounded by (EV9), or (IC2).
The preceding lemma supplies a subsequence of its Gauss traces
converging locally uniformly to a finite convex $g$.
The strong completed calculus gives the original coefficient
convergence $u_jS\to uS$, before applying any tropical test.
Consequently, for every $\varphi\in C_c^\infty(Q)$,
\[
 \begin{split}
 \langle t_*D(uS),\varphi\rangle
   &=\lim_j\langle u_jS,D(\varphi(t))\rangle\\
   &=\lim_j\int_Qg_j\,\Delta\varphi\,dy
     =\int_Qg\,\Delta\varphi\,dy
     =\langle\Delta g,\varphi\rangle.
 \end{split}
 \tag{SO6}
\]
The pullback $\varphi(t)$ is an original compact smooth test:
inverse images of compact real subboxes in this original torus
are compact polyannuli.  Each identity uses the already fixed
all-original calibration.  It does not define the completed
current by a selected toric test subfamily.

The positive measure $\Delta g$ has compact support in $Q$.
The exact frozen R10 convex compact-marginal rigidity theorem,
valid here for $n\geq2$, gives $\Delta g=0$.
Hence $t_*D(uS)=0$.  Positivity makes the original measure zero:
for every compact ordinary $L\subset U$, its mass is at most
the mass of $t^{-1}(t(L))$, which is the zero pushforward
mass of the compact set $t(L)$.  Radon inner regularity
then gives (SO5).

The datum (DO9) is finite, nonzero, compactly supported,
dominated by one half of an actual normalized tester measure,
and has finite energy-dual size, as already proved there.
It contradicts (SO5) if it were a completed source.
No claim is made that this special background exhausts arbitrary
backgrounds; one allowed background suffices to disprove the
unqualified universal source assertion.
\end{proof}

\paragraph{The resulting minimal equation gate.}
The strong completion solves a precise nonlinear Cauchy problem
and retains all-original closed curvature currents.  It does
not contain every solution of the larger relative linear
energy problem even when $q=1$.  The surviving convex trace
is an explicit obstruction.  A prospective general nonlinear
Dirichlet theorem must therefore specify an admissible operator
image or a different potential class; neither completion alone
nor a finite datum energy norm supplies the missing equation.
The analytic PS/PL relation quotient and low-real-degree CTO
are separate and are not used in this argument.
